\section{课题来源及研究的目的和意义}
\subsection{课题来源}
本课题来源于“复杂多层次异构微服务系统精准治理模型与关键技术”（国家
自然科学基金面上项目 62372140）。
\subsection{课题的研究背景}
近年来，以无人机（UAV）、电动垂直起降飞行器（eVTOL）等为代表的低空经济，
正作为战略性新兴产业，获得国家层面的高度重视与政策支持。
相关政策的密集出台，旨在推动构建现代化的空中交通和信息网络，
催生了物流配送、应急通信、环境监测、农林植保等海量新兴应用场景\cite{dikongjingji}。
这为构建一个全新的、运行于低空空域的服务网络生态提供了前所未有的发展机遇。

低空经济的核心价值在于“服务”，而非仅仅是飞行器本身。
上述各类应用场景，其本质都是通过部署在低空的、具备感知、计算和通信能力的移动节点来提供服务。
服务计算（Service Computing）作为研究服务系统建模、组合、交付与管理的学科，
为理解和构建这一新型服务生态系统提供了关键的理论框架和技术支撑\cite{service-project}。
它指导我们如何将低空中的物理实体（如飞行器、传感器）抽象为可被发现、调用和组合的服务单元，
并对这些服务的全生命周期进行有效管理。

然而，将服务计算的理念应用于低空经济的实践，也引出了一系列全新的系统性挑战。
首先，服务环境的极端动态性是其最显著特征。
传统的服务系统大多构建于稳定的网络基础设施之上，而低空服务系统则依赖于高速移动、
资源受限的边缘节点构成的自组织网络。节点的机动性导致服务网络的拓扑结构瞬息万变，
服务间的可达性与链路质量呈现出高度不确定性\cite{qiu2023integrated}。
其次，服务交付的性能约束对系统提出了严苛要求。用户对微服务的请求需要得到及时响应，
同时服务的执行过程必须保持连续和稳定，这对系统的服务质量（QoS）保障机制构成了严峻考验。
\cite{drlrr}。

如在应急通信的场景下，高动态移动边缘设备因其灵活性、适应性、稳定性、多功能性和低维护成本
被广泛应用于灾后行动\cite{after-disaster-operation}、
野外火灾探测和监控\cite{wildfire-detection-monitoring}和
搜救行动\cite{search-and-rescue}等。
在地震、海啸等重大自然灾害之后，地面通信基础设施常被严重损毁，
高动态边缘设备能够作为一种灵活、抗毁的空中通信中继网络被迅速部署至灾区，
为救援队和指挥中心提供关键的应急服务\cite{after-disaster-operation}。
在此类场景下，如图\ref{fig:uav}左图所示，移动边缘设备的核心任务是扮演动态的空中通信基站角色，通过多跳路由的方式，
将灾区现场采集到的监控视频、传感器读数等重要数据，可靠且高效地中继回远程数据中心。

\begin{figure}[htbp]
\includegraphics[width=\textwidth]{uav.pdf}
\caption{高动态边缘服务系统协同工作和路由策略对比示意图}
\label{fig:uav}
\end{figure}

在这种极端动态和服务严苛的双重压力下，一个根本性的矛盾贯穿系统设计的始终：
如何平衡“服务请求的即时满足”与“整个系统长期的可持续运营”？
具体而言，当一个服务请求到达时，系统若采用贪心策略，
调度最近的节点进行响应以追求最低延迟，可能会导致该节点因能耗过载或脱离主网络而过早失效，
进而损害系统的整体服务能力和任务周期。
反之，若完全从系统长期能效的角度进行全局调度，则可能牺牲单次服务的响应性能。
这种“局部最优”与“全局最优”、“短期利益”与“长期目标”之间的冲突，
是高动态服务系统治理中的核心难题，也是现有技术方案普遍面临的瓶颈\cite{1025218244.nh}。

解决这一核心矛盾，需要对支撑服务交付的底层机制进行深入研究。
一方面，移动服务节点在物理空间的部署形态，即物理服务拓扑，
直接决定了服务与用户间的距离以及网络连通性的基础。
另一方面，在动态变化的拓扑之上，服务数据在节点间的路由转发策略，如图\ref{fig:uav}右图所示，
则决定了能量和负载在全网的分配，影响着系统的生存周期\cite{routing-strategy}。
现有研究往往将拓扑控制与路由决策割裂处理，忽略了两者间的强耦合性：
一个节点的“最优位置”取决于数据能否通过高效的路径传达；
反之，一条“最优路由”的存在，又依赖于节点处在合适的物理位置。
这种解耦优化的方式难以实现系统级的全局最优。

为了系统性地应对上述挑战，本课题将采用一种演进式的研究思路，
旨在构建一个能够协同优化拓扑与路由的智能治理框架。
为此，本文提出并聚焦于以下三个核心研究问题：

\begin{enumerate}
\item \textbf{面向高动态服务系统的自适应拓扑优化方法研究}：旨在设计一种自适应的拓扑控制策略，使移动服务节点集群能够根据宏观的服务需求分布和网络连通性要求，动态调整其物理部署形态，为上层服务的高效交付构建稳健的基础。
\item \textbf{面向高动态服务系统的分布式能量均衡路由策略研究}：旨在基于李雅普诺夫优化理论，设计一种分布式的路由决策机制，在保障数据传输性能的同时，实现网络整体能耗的均衡，避免因“能量热点”导致关键节点过早失效，从而最大化系统的服务周期。
\item \textbf{拓扑控制与路由决策的联合优化方法研究}：在前两个问题的基础上，旨在构建一个反馈驱动的联合优化框架，将拓扑调整的长期收益与路由决策的短期效用进行统一建模和协同求解，以实现系统在“即时性”与“可持续性”这对核心矛盾间的动态最优平衡。
\end{enumerate}

\subsection{课题的研究意义}
本课题旨在系统性地解决高动态服务网络中的治理难题，其理论与实践意义紧密围绕所提出的研究问题展开。

在理论层面，本研究期望构建一个系统性的优化框架，
以应对现有研究在处理多重冲突目标（如性能与能耗）时缺乏系统性理论指导的不足。
为此，本课题将引入并扩展李雅普诺夫优化理论，
通过将系统的长期稳定运行目标分解为每个节点在每个时刻的本地化决策，
为解决高动态环境下的性能与能耗矛盾提供一种具备理论性能保证的优化范式。
同时，针对当前研究大多将拓扑控制与路由决策解耦处理的现状，
本课题的最终目标是探索一种基于深度强化学习的联合优化框架。
该框架将拓扑调整的长期收益与路由决策的短期效用进行统一建模和协同求解，
为解决复杂网络系统中的多目标、强耦合问题提供了新的理论视角。
此外，多数现有路由策略的决策依据是瞬时和局部的，缺乏对网络长期状态和生存性的考量。
本研究将“能量均衡”和“系统生命周期”作为核心优化目标，
旨在弥补当前在高动态网络长期规划与生存性保障方面的研究空白。

在实践层面，本研究成果可为低空经济背景下的新兴应用提供关键技术储备。
一个能够自主调节拓扑并智能规划路由的无人机网络，将极大提升其在城市巡检、
物流配送等任务中的成功率和持续时间，为相关产业的发展提供坚实的技术支撑。
同时，本课题将服务计算的宏观理念与底层网络技术的具体实现相结合，为服务计算理论在高动态、
资源受限等挑战性环境下的落地应用提供了具体的算法解决方案和可行性路径。
最终，通过引入联合优化机制，本研究旨在使高动态网络从一个被动执行指令的工具集，
演进为一个能够根据自身状态和外部需求进行自我感知、自我调整的智能有机体，
这对于提升未来自主系统集群的智能化与协同作业能力具有重要意义。
\section{国内外在该方向的研究现状及分析}
本章节将对国内外在动态网络管理领域的相关研究现状进行综述与分析。
论述将围绕本课题提出的核心研究方向展开，
首先梳理并分析高动态服务系统中的物理拓扑控制方法，
聚焦于如何通过调整节点位置来为服务交付构建物理基础；
其次，将探讨服务系统中的能量均衡路由策略，重点关注在高动态、
能量受限环境下的数据转发决策问题；最后，将对当前国内外相关领域研究的共性不足进行归纳，
明确本研究的切入点与创新方向。
\subsection{高动态服务系统拓扑控制研究现状}
在高动态服务系统中，移动节点物理位置的动态调整，即拓扑控制，
是保障服务覆盖与网络连通性的基础。现有研究主要集中于节点部署优化、
网络组织结构（即集群化）以及动态拓扑维护等方面。

在节点部署优化方面，一个主要的研究方向是将移动节点（尤其是无人机）作为空中基站
（Aerial Base Stations, ABS）或中继节点，通过优化其三维空间位置，
为地面用户提供最大化的服务覆盖范围或最优的通信质量\cite{ali2020uav}。
部分研究考虑了空中基站与地面基站（TBS）共存的蜂窝网络，
重点分析了通信干扰下的最优部署策略\cite{zhang2019iot}。在地面基础设施被破坏的场景下，
动态部署空中基站以提供紧急无线覆盖成为研究热点，
通常会引入移动模型对用户行为进行仿真\cite{barick2025uav}。

为了对大规模无序节点进行有效组织和管理，集群化是拓扑控制中的一种关键技术。
传统的集群算法被广泛应用于移动自组网中，例如基于最小身份ID的算法、
基于最大连通度的算法以及基于移动性预测的集群协议（MPBC）\cite{ni2011mpbc}。
其中，加权集群算法（WCA）因其能够综合考虑节点度、移动性、剩余能量等多个指标，
对多变场景的适应性而备受关注\cite{chatterjee2002wca}。近年来，随着数据驱动方法的发展，
研究者开始运用机器学习算法进行网络组织。例如，利用K-Means\cite{hartigan1979algorithm}或K-Mediods\cite{rdusseeun1987clustering}
等聚类技术对地面服务节点或用户进行分组，再基于分组结果规划中继节点的部署与轨迹。

在高动态环境中，拓扑的建立与维护是一个持续的挑战。
飞行自组织网络（FANET）因其高动态特性，成为拓扑维护研究的典型场景。
为应对节点高速移动导致的拓扑快速变化，研究者提出了多种动态拓扑维护算法。
例如，基于粒子群优化（PSO）的拓扑构建算法旨在保障端到端的通信质量\cite{kim2017topology}。
同时，为了应对时变动态的网络环境，联合优化任务分配和拓扑管理的问题也得到了关注\cite{bekmezci2014connected}。
为保证网络的稳定性，需要不断维护拓扑结构\cite{kim2019joint}，已有研究基于梯度下降增量调整拓扑，
或利用图的连通性，提出基于位置运动和深度强化学习的拓扑控制算法\cite{kim2018integrated}。
\subsection{高动态服务系统能量均衡路由研究现状}
在单个集群或扁平化网络架构内，如何设计高效的路由策略是保障高动态服务系统提供高效稳定服务的核心。现有研究根据其技术路径，大致可分为传统地理路由、启发式算法、强化学习以及李雅普诺夫优化等几类。

传统路由协议，尤其是地理位置路由，因其简洁性在早期自组网研究中得到了广泛应用。
最具代表性的GPSR\cite{gpsr}协议采用一种基于位置的贪心路由策略，即节点总是选择地理位置上最接近目的地的邻居作为下一跳。
OLSR\cite{olsr}协议作为一种典型的链路状态协议，它的核心思想是让网络中的每个节点都拥有完整的网络拓扑信息，从而能够独立计算出到达任意目的节点的最佳路径。

为克服传统方法的局限性，在复杂的网络环境中寻找最优路径，一些研究探索了启发式算法的应用。
例如，在水下环境中，有研究采用蚁群优化算法（ACO）\cite{aco}来迭代收敛到最优路径。
这类方法这类方法通过模拟蚂蚁探索路径并留下信息素，依据信息价值等指标来评估路径质量，从而进行动态的数据包路由。
也有研究者利用AFS算法\cite{afsa}来解决应用层的多播路由问题。这些生物启发式算法在特定场景下表现出良好的寻优能力。

近年来，以强化学习（RL）为代表的人工智能技术因其强大的环境自适应能力，在高动态网络路由领域受到了极大关注。
这类方法将路由选择过程视为一个马尔可夫决策过程（MDP），使无人机能够通过与环境的交互和试错，自主学习最优的转发策略。
QTAR\cite{qtar}提出利用两跳邻居信息来扩展节点的局部网络拓扑视野，从而做出更可靠的路由决策。
其设计的奖励函数综合考虑了邻居节点的位置、速度、延迟和剩余能量等多重指标。
QGeo\cite{qgeo}则将Q-learning与地理路由相结合，节点位置信息增强了路由发现过程，应用Q学习来选择高移动性无人机节点之间的最优路由。
为处理高维连续的状态空间，基于深度强化学习（DRL）的策略也已被引入：
DRLRR\cite{drlrr}是一种为城市应急通信网络设计的弹性路由策略，它将路由决策过程建模为MDP，并采用双深度Q网络（DDQN）算法进行训练，以适应链路中断等网络拓扑的剧烈波动。
而DRLFR\cite{drlfr}利用DRL来处理高机动性无人机带来的延迟和信道增益的状态量化误差，从而协同决定转发行为和中继功率。

针对高动态移动节点对于能耗的敏感性，对其进行长期的优化非常重要。
李雅普诺夫优化框架能够将一个长期的、跨时域的优化问题，分解为一系列在每个独立时隙内求解的、确定性的低复杂度问题。
一部分研究聚焦于利用李雅普诺夫优化来保障数据队列的稳定性以最小化延迟。
例如，通过引入虚拟队列来对路由施加跳数约束\cite{related1}，或将其与DRL结合以自适应调整路由策略\cite{related2}。
另一类工作运用李雅普诺夫框架对无人机的能量进行显式控制，
有研究提出了一种基于李雅普诺夫的方法，用于联合优化太阳能无人机的三维轨迹和资源分配\cite{related3}，目标是在保障能量供给稳定的前提下，最大化系统吞吐量。
在计算卸载场景中，Lin等人\cite{related4}设计了虚拟电池能量队列来构建李雅普诺夫函数，从而为作为计算核心的无人机集群设计了能够维持各节点长期能量稳定的卸载算法。
% \subsection{多集群分层路由研究现状}
% 当高动态边缘服务系统规模增大时，为保证其可扩展性，采用分层路由架构成为主流的技术路径。
% 该架构的核心在于将网络划分为若干集群，并通过簇头进行分级管理与路由。
% 簇头不仅负责管理簇内通信，还与其他簇头共同构成一个骨干网络，以处理跨集群的数据转发。
% 现有研究主要围绕如何高效地构建和维护集群，以及如何在此基础上进行路由展开，这包含了三个关键环节：簇头选举、集群维护以及簇间路由。
% 如何选举出合适的簇头、在拓扑剧变中动态维护集群结构，以及如何实现高效的跨集群通信，是决定整个大规模网络性能的关键，现有的研究可能将其中某些环节进行合并。

% Hosseinzadeh等人提出了一种基于鲸鱼优化算法的飞行自组织网络（FANET）智能聚类方案\cite{woa}，
% 通过模仿鲸鱼的觅食行为来选择最佳簇头，允许基于无人机移动性和网络条件动态调整簇头，从而提高网络稳定性和降低能耗。
% 通过选择能量储备更高、连接性更好的簇头，该方法最大限度地减少了重新聚类的需要，减少了路由开销并增强了网络寿命。
% IKAD-FANET\cite{ikad-fanet}提出了一种两阶段选举方法：首先，通过改进的K-Means算法根据空间邻近性选出初始簇头；
% 然后，再利用密度峰值聚类（DPC）算法，根据节点密度对初始簇头进行二次筛选和优化，以确保最终当选的簇头位于节点密集区域，从而提升集群的稳定性。
% Song\cite{song}等人通过泛洪路由利用聚类方法和泊松集群过程来划分无人机网络，从而最小化跳数，降低延迟违规概率。

% ICRA\cite{icra}提出了一个完整的框架，其核心是利用强化学习模型来动态调整聚类策略本身。
% 在簇头选举阶段，其将节点的效用值定义为四个因素的加权和：剩余能量、节点中心度、与邻居的速度相似性以及与邻居的链路保持时间。
% 它利用强化学习模型，根据网络状态的反馈（如拓扑稳定性和能耗变化）来动态地学习和调整这些因素的权重，而不是采用固定的权重值。
% 在集群维护方面，当部署在地面站的强化学习模型学习到新的、更优的聚类策略（即权重组合）时，
% 会向全网广播该策略，所有边缘设备在接收到新策略后将触发一次全局的重聚类过程，以适应宏观网络状态的变化。
% 在簇间路由过程中，ICRA采用基于混合骨干网的方式，除了簇头节点外，那些位于集群边缘、且能够与相邻集群节点通信的普通无人机，
% 可以申请成为“簇间转发节点”，以增强网络的灵活性。

% K-MORP\cite{k-morp}主要是做了动态集群维护和簇间路由的工作，
% 当一个无人机成员移动并超出了其簇头预定义的距离阈值时，其会触发网络的重聚类过程。当检测到链路断开时，新的簇头会发起路径重发现过程。
% 当源节点需要向其他集群的节点发送数据时，它首先将数据包发送给本集群的簇头，随后通过簇头间的骨干网络，将数据包转发至目的节点所在集群的簇头。
% 在簇间路由上，Yang\cite{caconet}等人提出了一种基于混合蚁群优化的FANET集群间路由协议，利用ACO的仿生方法来寻找最优集群间路由，同时最小化开销，实现了对快速网络拓扑变化的高适应性。
\subsection{现有工作的不足}
尽管现有研究为高动态服务系统网络治理提供了多种解决方案，
但无论是在物理拓扑控制还是在数据路由策略上，都存在着制约系统长期、高效运行的共性不足。

在高动态边缘服务系统中，现有拓扑控制方法的不足主要体现在其决策的静态性和对服务动态需求的感知缺乏上。
多数基于静态优化的节点部署方法，如将无人机视为固定空中基站，
其优化目标往往基于对用户或服务需求的快照式假设，缺乏对环境动态变化的适应能力，
一旦服务需求发生迁移，其拓扑便迅速劣化为次优。基于集群化的拓扑组织方法，
虽然提升了网络的可管理性，但其效用函数中的权重（如能量、连接度等）往往是预先设定的，
使其难以根据任务不同阶段的演进动态调整拓扑的侧重点。即便是引入了强化学习等自适应方法，
其奖励函数的设计也通常是短视的，侧重于最大化即时的拓扑指标（如覆盖率、连通度），
而未能有效量化一个拓扑形态对未来数据路由效率和系统整体能耗的长期影响。
因此，现有拓扑控制方法本质上是“为拓扑而拓扑”，缺乏一种面向长期服务效能的、
具有前瞻性的战略规划能力。

而在高动态边缘服务网络的路由问题上，现有路由策略的不足主要体现在其决策的短视性和对长期网络状态的忽略上。
大多数传统路由策略，如基于地理位置的贪心路由，其决策依据是瞬时和局部的最优性，
总是倾向于选择处于某些关键位置的节点进行中继，导致这些节点被过度使用，能量被过早耗尽。
节点的过早失效会直接引发网络分裂和服务中断，这对于需要长时间提供稳定的应急通信服务是致命的。
启发式算法虽然具备寻优能力，但其高计算复杂度和迭代收敛过程使其难以适应高动态网络拓扑高速变化下的实时决策需求。
强化学习算法具备自适应能力，但面临收敛速度慢、在动态环境中需要高昂代价再训练等问题，
更重要的是，它们通常缺乏一个明确的机制来对长期能耗进行管理，从而无法为服务的连续性提供硬性保障。
尽管李雅普诺夫优化是解决长期随机优化问题的有力工具，但现有相关工作要么仅关注数据队列稳定以降低延迟，
而忽略了能量管理；要么在非路由场景（如能量收集的轨迹规划）中进行能量控制，
而没有解决分布式、多跳路由场景下的能量均衡问题。
因此，在利用该理论解决路由中的长期能量健康问题上，存在明显的研究空白。

更关键的一点在于，现有工作普遍将拓扑控制与路由决策解耦处理。
拓扑控制算法在规划节点位置时，对该决策将如何影响后续的路由选择是“盲目”的；
而路由算法则只能在给定的、未必合理的拓扑约束下被动地寻找“最优”路径。
这种研究范式的割裂，从根本上限制了系统所能达到的全局性能上限。
一个在拓扑层面看似优秀的形态，可能在路由层面造成了难以解决的拥塞瓶颈。
这是当前高动态服务系统治理面临的核心挑战，也正是本研究试图通过联合优化来突破的关键所在。

% 在面向大规模网络的多集群架构中，现有工作虽然通过集群化解决了部分扩展性问题，
% 但在集群的构建、维护和路由的长期规划上，依然存在显著不足。
% 静态、启发式的集群构建策略依赖一个预设的、静态的效用函数来选举簇头，这种方法的根本缺陷在于其缺乏长期自适应能力。
% 随着网络整体能量状态和拓扑结构的变化，可能在中后期变得不再适用，从而选出不合适的簇头，导致集群结构频繁重构，稳定性差。
% 而即便是引入了强化学习等自适应方法，其决策目标也往往存在偏差。
% 它们通常将问题建模为最大化长期累积折扣奖励，这使得策略倾向于瞬时最优，而缺乏对任务执行时间窗口的明确感知，
% 无法解决“如何在确保任务周期内能量充足的前提下，最大化系统性能”这一核心问题。
% 在集群维护方面，现有的集群维护大多是被动和反应式的，例如，仅当簇头能量低于某个阈值或与成员失联后，才触发更新或重选。
% 这种模式缺乏前瞻性的长期规划，未能将集群的“能量健康”作为一个连续状态进行主动管理。
% 在簇间路由方面，尽管现有工作设计了簇头间的路由机制，但这些机制在选择下一跳簇头时，
% 往往再次陷入了地理位置最优等短视的决策模式，
% 这种做法会过度消耗某些关键网关簇头的能量，同样会导致骨干网的拓扑割裂，造成大片网络区域的通信中断。
\section{主要研究内容及研究方案}
\subsection{主要研究内容}
\begin{figure}[htbp]
\includegraphics[width=\textwidth]{research.pdf}
\caption{研究方案框架图}
\label{fig:research}
\end{figure}
图\ref{fig:research}展示了本课题的主要研究内容框架，共有三个研究问题：
\begin{enumerate}
  \item \textbf{面向高动态服务系统的自适应拓扑优化方法研究}：
本部分内容旨在解决移动服务节点集群的物理组织与形态管理问题。
在高动态服务系统中，一个优化的物理拓扑是保障服务覆盖、降低通信延迟和维持网络连通性的基础。
本研究将以无人机集群为典型场景，探索在满足服务需求和网络稳定双重约束下，
如何通过分布式协同决策，动态优化每个服务节点的位置。
研究将具体围绕三个核心挑战展开：如何将无序的节点自适应地组织成高效的集群结构；
在集群内部，如何协同调整各节点位置以最优地响应时变的服务需求；
以及在确定目标位置后，如何为多个节点规划安全且保障通信的运动轨迹。
  \item \textbf{面向高动态服务系统的分布式能量均衡路由策略研究}：
本部分内容聚焦于在动态变化的拓扑基础上，解决服务数据的路由转发决策问题。
其核心挑战在于平衡网络传输性能（如吞吐量）与系统长期生存性（如能量消耗）之间的内在矛盾。
在高动态、能量受限的环境下，短视的、仅追求最高性能的路由决策极易导致关键节点能量过早耗尽，
从而引发网络分裂。本研究将引入并应用李雅普诺夫优化理论，
旨在将长期的、随机的系统优化目标，
转化为每个节点在每个时刻的、仅依赖局部信息的确定性优化问题，
从而设计一种分布式的、能够在理论上实现能量均衡的路由与功率联合决策方案。
  \item \textbf{拓扑控制与路由决策的联合优化方法研究}：
本部分旨在克服前两个问题解耦优化所带来的局限性。拓扑与路由的决策是强耦合的，
一个优化的拓扑形态应为高效路由创造条件，而高效的路由策略也反过来依赖于合理的拓扑布局。
本研究旨在构建一个统一的智能决策框架，将整个服务集群视为一个多智能体系统，
利用深度强化学习方法，学习一个能够协同调整拓扑和路由的端到端联合治理策略。
该策略的目标将不再是孤立的拓扑或路由指标，而是旨在最大化一个能够反映系统长期、
综合服务效能的奖励函数，从而在“即时性”与“可持续性”的矛盾中找到动态最优解。
\end{enumerate}
\subsection{研究方案}
\subsubsection{面向高动态服务系统的自适应拓扑优化方法研究}
本方案旨在设计一套包含集群构建、簇内位置协同和轨迹规划三个环节的自适应拓扑优化流程。
在研究方法上，首先将采用基于模糊C均值（FCM）的动态分簇算法，通过将节点的地理位置、
剩余能量等多维信息作为特征向量，实现节点的自适应集群化，
其“模糊”隶属度特性可有效应对节点在簇边缘的平滑过渡。接着，在集群内部，
将采用基于李雅普诺夫优化的拓扑效用最大化方法，通过构建衡量服务积压和机动能耗的虚拟队列，
并最小化单步“漂移+惩罚”项，来推导出指导节点移动的效用函数，使集群形态趋向长期效用最优。
最后，在确定节点目标位置后，将采用基于通信感知的冲突搜索（CBS）算法，
通过将“通信中断”引入为一种新的冲突类型，为多节点生成兼顾物理安全与通信稳定的协同轨迹。

为验证本方案的有效性，将通过仿真实验，与多种基线算法进行横向对比，包括静态队形策略、
贪婪服务追踪策略以及不考虑通信约束的基础强化学习轨迹规划算法。
评估将围绕多个关键指标展开，如服务请求覆盖率、平均网络连接时间、
拓扑变化率、总机动能耗以及服务请求平均满足延迟。
\subsubsection{面向高动态服务系统的分布式能量均衡路由策略研究}
本方案旨在构建一套基于李雅普诺夫优化理论的分布式路由与功率联合决策方案。
在研究方法上，将运用漂移加惩罚理论，通过为每个节点定义物理数据队列和能量赤字虚拟队列，
将最大化长期吞吐量且满足长期能量约束的全局问题，
转化为最小化每个节点在单时隙内“漂移+惩罚”项的本地化问题。
该本地化问题进一步分解为自适应传输功率控制和下一跳转发决策两个子问题。
前者通过对单步效用函数求解得到最优传输功率；后者则通过构建一个融合了理论优化收益、
节点能量与地理进度的复合链路权重来选择最优下一跳。为保障决策信息的时效性，
还将设计一种基于移动预测的自适应HELLO消息广播间隔调整机制，
以在信息新鲜度与网络开销间取得平衡。

为进行性能评估，本方案将与多种经典及先进的路由协议进行对比，
主要包括经典的地理路由策略GPSR以及其复杂均衡策略LD-GPSR、
基于蚁群算法的路由策略ACO和基于深度强化学习的DRLFR以及作为理论基础参照的LEADR策略。
对比将通过一系列关键性能指标进行量化，包括网络总吞吐量、平均数据包端到端延迟、
作为核心均衡性度量的全网节点剩余能量的方差，以及衡量系统生存能力的网络生命周期。
\subsubsection{拓扑控制与路由决策的联合优化方法研究}
本方案旨在采用多智能体深度强化学习方法，构建一个端到端的联合治理模型。
在研究方法上，将采用“中心化训练，去中心化执行”的学习范式，
并选用如MADDPG等先进算法进行训练。为实现联合优化，将设计一个融合了节点自身、
局部拓扑与宏观服务信息的联合状态空间，
以及一个包含离散移动与路由决策的二维元组式联合动作空间，
迫使智能体在单步决策中对两者进行协同考量。为引导智能体学习到长期最优策略，
将设计一个融合了即时任务收益与长期生存代价的量化奖励函数，
该函数将对成功的数据转发和任务靠近给予正奖励，对能量过度消耗、链路中断等行为给予负奖励。

为证明联合优化相比于解耦优化的优越性，将进行纵向消融实验。
实验将对比本方案的联合治理模型、消融了智能拓扑控制的模型、消融了智能路由策略的模型，
以及完全由简单基线算法构成的解耦模型。评估将采用系统级的综合指标，
以衡量长期效能，包括任务总完成率、平均服务响应时间、系统总能效以及系统生命周期，
旨在定量地展现本研究提出的联合治理框架所带来的协同增益效应。
\section{预期达到的目标}
以下将概述针对三个研究问题预期达到的目标：
\begin{enumerate}
  \item \textbf{提出一套自适应的物理拓扑优化方法，为高动态服务交付构建基础}：预期将设计并实现一套能够根据服务需求和节点状态，动态组织和管理集群拓扑的算法。该方法应能实现无序节点的自适应集群化，并在簇内自主调整节点位置，以动态平衡服务覆盖范围与内部通信稳定性。通过仿真实验，预期该方法在服务请求覆盖率、网络连通性维持时间等关键指标上，将定量地优于静态队形或贪婪服务追踪等基线策略。
  \item \textbf{设计一套分布式的能量均衡路由策略，显著提升系统的长期生存能力}：预期将构建并验证一个基于李雅普诺夫优化理论的分布式路由算法。该算法能够在每个节点上，仅利用局部信息做出最优的下一跳转发和传输功率的联合决策。核心目标是有效均衡网络整体的能量消耗，避免“能量热点”的产生。通过仿真实验，预期该策略在显著延长网络生命周期（以首个节点失效时间衡量）和提升网络能量均衡度（以节点剩余能量方差衡量）的同时，依然能保持与现有先进路由协议（如GPSR、DRLFR）相媲美的网络吞吐量和数据包投递率。
  \item \textbf{构建一个拓扑与路由的联合治理框架，验证协同优化的系统性增益}：预期将设计并实现一个基于多智能体深度强化学习的联合优化框架。该框架能够学习到一个端到端的智能策略，在单步决策中协同完成移动和路由的选择。最终目标是使系统能够在“即时服务性能”与“长期系统生存”这对核心矛盾中做出动态的、接近最优的权衡。通过与解耦优化方案（即将前两个研究问题中优化的算法进行简单组合）进行纵向消融实验对比，预期本联合治理框架在系统级的综合评估指标（如任务总完成率、系统总能效）上展现出协同增益效应，从而验证联合优化的必要性与优越性。
\end{enumerate}
\section{已完成的研究工作与进度安排}
\subsection{已完成的工作}
目前，本课题已围绕核心研究问题开展了前期的理论分析与问题建模工作，
为后续的算法设计、系统实现与实验验证奠定了基础。
已完成的工作主要集中于对高动态服务系统的物理拓扑优化和能量均衡路由这两个核心子问题进行数学形式化描述。
% \subsubsection{云边端复杂软件系统运维和运行平台扩展架构设计}
% 本课题的研究将基于现有的云边端复杂软件系统运维和运行平台进行二次开发与功能扩展。
% 该平台架构如图\ref{fig:platform_arch}左侧所示，底层基于Kubernetes，已初步具备纳管异构集群、应用部署和基础监控的能力。
% 在此基础上，为应对高动态移动边缘网络的独特挑战，我们设计了如图\ref{fig:platform_arch}右侧所示的扩展架构，
% 他们是当前平台三个模块的架构设计与具体细化，三个模块协同工作，共同为高动态网络提供完整的平台级支持：
% \begin{figure}[htbp]
% \centering
% \includegraphics[width=\textwidth]{arch.pdf}
% \caption{面向云边端复杂软件系统运维和运行平台扩展架构示意图}
% \label{fig:platform_arch}
% \end{figure}

% \begin{itemize}
% \item \textbf{智能运维与算法运行引擎}：位于架构顶层，它负责全面感知系统状态，运行各类优化算法（如路由策略），并根据算法结果生成运维决策，实现智能化的闭环控制。
% \item \textbf{调度引擎}：作为平台的核心枢纽，负责资源的规划与任务的分配。它根据上层运维引擎的决策或用户的直接请求，将应用和服务智能地调度到最合适的物理节点或集群。
% \item \textbf{网络互联及路由配置引擎}：作为平台的底层网络基础，负责打通不同集群之间的网络壁垒，并提供动态配置路由规则的能力，是实现大规模、跨集群服务协同的前提。
% \end{itemize}

% 本人的研究工作将聚焦于\textbf{智能运维与算法运行引擎}的设计与实现。
% 该引擎是连接底层状态感知与上层智能决策的关键桥梁，旨在为平台赋予“看得见、看得懂、能决策、能执行”的核心能力。
% 其具体实现方案将围绕数据采集、状态监控和算法编排三个核心模块展开。
% \textbf{1. 数据采集与管理模块设计}
% 该模块是智能运维的数据基石，负责从多源异构的底层系统中全面、实时地采集运维数据。
% 在功能上，它将周期性地采集来自集群、网络和应用等多个维度的时序指标；
% 集中聚合各底层组件的日志；通过事件总线监听并汇聚来自各类系统的离散事件；
% 并定期扫描集群以自动发现服务间的关联关系，构建动态的网络拓扑。
% 在技术选型上，指标采集将采用业界主流的\textbf{Prometheus}，
% 日志聚合选用\textbf{Grafana Loki}方案，
% 而事件与拓扑数据将通过调用Kubernetes API等接口收集，由基于\textbf{SpringBoot}的后端服务整合并持久化存储。
% \textbf{2. 核心状态监控与可视化模块设计}
% 该模块负责将采集到的海量、多维数据转化为直观、可交互的可视化视图，
% 为运维人员和上层算法提供决策依据。功能上，它不仅提供标准的集群资源监控视图，
% 还将针对高动态边缘节点，扩展展示其三维地理位置、剩余电量、数据队列积压等特有状态。
% 同时，它会以图形化方式动态展示服务间的调用关系与关键网络指标，
% 并为具体应用提供性能观测仪表盘。技术选型上，将采用\textbf{Grafana}作为主要的可视化工具，
% 并借助\textbf{Istio}服务网格与\textbf{Jaeger}实现网络流量与链路追踪。
% 对于动态网络拓扑等高度定制化的视图，将由\textbf{Vue}前端框架开发实现。
% \textbf{3. 算法编排与管理模块设计}
% 该模块是连接“分析”与“计划/执行”的关键，负责将本课题研究的路由优化等算法模型，
% 以标准化的方式在平台上运行、管理和调度。功能上，它允许用户通过界面发起分析任务，
% 灵活配置数据源与算法参数，并对算法任务的执行全流程进行追踪，最终将算法的输出结果以可视化方式进行展示和对比。
% 技术选型上，算法的调度与生命周期管理将基于云原生工作流引擎\textbf{Argo Workflows}实现，
% 其执行日志将输出至\textbf{Loki}并通过\textbf{Grafana}进行实时展示
\subsubsection{高动态服务系统物理拓扑优化问题建模}
本部分工作对单个服务集群的物理与通信环境进行了形式化建模，旨在为自适应拓扑控制策略的研究提供数学基础。

我们考虑一个由地面用户、高移动服务节点和服务集群组成的多层次系统，其集合定义如下：
\begin{itemize}
	\item \textbf{地面用户集合}：$\mathcal{X}_{U} = \{\vec{u}_1, \vec{u}_2, \ldots, \vec{u}_l\}$，其中 $\vec{u}_i = (x_i, y_i, z_i)$ 表示第 $i$ 个用户的三维位置坐标。
	\item \textbf{无人机集合}：$\mathcal{X}_{V} = \{\vec{v}_1, \vec{v}_2, \ldots, \vec{v}_n\}$，其中 $\vec{v}_j = (x_j, y_j, z_j, \varphi_j, \alpha_j, F_j, M_j)$ 表示第 $j$ 个服务节点的位置、水平飞行角度、垂直飞行角度、飞行速度，以及CPU频率 $F_j$ （周期/秒）和内存大小 $M_j$（比特）。
	\item \textbf{任务集合}：$\mathcal{X}_{T} = \{\vec{t}_1, \vec{t}_2, \ldots, \vec{t}_n\}$，其中 $\vec{t}_k = (u_{k\_src},  d_k)$ 表示一个任务需求，由源用户 $u_{k\_src}$ 发出，数据量为 $d_k$ （比特）。
	\item \textbf{无人机集群集合：} $\mathcal{X}_{C} = \{C_1, C_2, \ldots, C_K\}$，其中 $K$ 表示集群的总数量，$C_k$ 是第 $k$ 个无人机集群。每个集群 $C_k$ 都是无人机集合 $\mathcal{X}_{V}$ 的一个子集，即 $C_k \subseteq \mathcal{X}_{V}$，并且所有集群的并集覆盖了所有无人机，即 $\bigcup_{k=1}^{K} C_k = \mathcal{X}_{V}$。
	\item \textbf{任务-无人机分配序列：}我们定义一个任务分配映射 $\mathcal{M}_{DT}$，它描述了每个任务如何分配给一个或多个无人机。该映射可以用一个二元组集合表示，其形式为 $\{(\vec{t}_k, \mathbf{m}_k) \mid \vec{t}_k \in \mathcal{X}_T\}$。其中 $\mathbf{m}_k$ 是一个向量，其分量 $m_{jk}$ 表示为完成任务 $\vec{t}_k$ 而分配给无人机 $j$ 的计算量大小。该分配必须满足约束：$\sum_{j \in \mathcal{X}_{V}} m_{jk} = d_k$。这确保了整个任务的计算量被完全分配，并且可以由多个无人机协作完成。
\end{itemize}

假设用户使用TDMA将计算任务卸载到无人机上。用户与无人机间的交互通常被视为可能性LoS模型，无人机间（A2A）的传输通常视为视距传输，适用于自由空间模型。

对于无人机 $v_i$，其三维坐标为 $\omega_i(t) = [x_i(t), y_i(t), h_i(t)]^T$。假设其飞行距离为 $l_i(t)$，水平方向飞行角度为 $\varphi_i(t)$，则
\begin{align}
	x_i(t+1) &= x_i(t) + l_i(t)\cos(\varphi_i(t))  \\
	y_i(t+1) &= y_i(t) + l_i(t)\sin(\varphi_i(t)) 
\end{align}

假设有一个最大仰角 $\psi_i$，则最大地面水平辐射半径为
\begin{equation}
	C^{\max}_i(t) = h_i(t)\tan(\psi_i) 
\end{equation}

对于 $\forall u_i$，总是 $\exists v_i$，使得 $\delta(u_i, v_i) \leq C^{\max}_i(t)$，其中 $\delta(u, v)$ 用来指示两个节点的地理距离。

由于无人机有限的水平飞行和垂直飞行速度，无人机有飞行距离限制：
\begin{align}
	\begin{split}
	l_i(t) &= \sqrt{(x_i(t+1) - x_i(t))^2 + (y_i(t+1) - y_i(t))^2} \leq L^h_{\max}  \\
	\Delta h_i(t) &= |h_i(t+1) - h_i(t)| \leq L^v_{\max} 
	\end{split}
\end{align}

其中 $L^h_{\max}$ 和 $L^v_{\max}$ 分别是水平和竖直方向的最大距离。对于无人机的飞行范围，不能超出区域 S，同时考虑到可飞范围，对于飞行高度亦有限制，$H_{\min}$ 和 $H_{\max}$ 是无人机在竖直方向上的最小和最大高度：
\begin{equation}
	\begin{split}
		&0 \leqslant x_i(t) \leqslant X, 0 \leqslant y_i(t) \leqslant Y \\
		&H_{min} \leqslant h_i(t) \leqslant H_{max}
	\end{split}
\end{equation}

为了保证两个无人机处在安全距离，对于任意两个无人机 $u_i$ 和 $u_j$，不能近于最小距离
\begin{equation}
	\delta(v_i, v_j) = ||\omega_i(t) - \omega_j(t)|| \geq D_{\min} 
\end{equation}

同时，为保证覆盖率，对于每一个地面用户，在任意时刻，总存在一个无人机，二者处于通信范围中。即：
\begin{equation}
	\forall i \in U, \exists j ,\delta(u_i,v_j) \leqslant C^{\max}_i(t)
	\quad \quad
\end{equation}

本文不关注任务是如何分配给无人机和无人机集群的，这属于任务卸载的范畴。本文关注的是，在任务分配之后，无人机集群应该如何重组、无人机应该如何行动。

为了分析系统的动态行为，我们对任务的到达和执行进行建模。

任务生产量是指在每个时隙 $t$ 内，由地面用户新产生的任务总量。
系统的总任务生产量可表示为：
\begin{equation}
	\Psi^{gene} = \sum_{\vec{t}_k \in \mathcal{X}_{T} } d_k 
\end{equation}

一个任务可能需要多个时隙完成，一个时隙内可能无法完成一个任务。在一个时隙$\Delta$内，无人机 $j$ 在时隙 $t$ 内的计算能力为 $F_j \times \Delta$。计算密度为 $\rho_k$，单位为周期/比特，表示处理单位比特数据所需的CPU周期数。可以完成的任务量为
\begin{equation}
	\Psi^{up} = \frac{F_j \times \Delta}{\rho_k}
\end{equation}

已知每个任务 $k$ 分配给无人机 $j$ 的数据量 $s_{jk}$。那么无人机完成任务所需的时间为
\begin{equation}
	T_{jk}^{comp} = \max\{\Delta,\frac{s_{jk} \times \rho_k}{F_j}\}
\end{equation}

解决的任务量为
\begin{equation}
	\Psi_{j,k} = \min\{\Psi^{up}, s_{jk}\}
\end{equation}

任务解决量是指在每个时隙 $t$ 内，所有无人机共同完成的计算任务总量。这部分任务量取决于无人机群体的总算力。

整个无人机集群的总任务解决量为所有无人机解决量之和：
\begin{equation}
	\Psi^{sys} = \sum_{j \in \mathcal{X}_V} \Psi_{j}(t) = \sum_{j \in \mathcal{X}_V} \sum_{k \in \mathcal{X}_D} \Psi_{j,k}
\end{equation}

由于无人机作为边缘服务器，其为任务主导性。相对复杂计算任务花费的高能耗，通信和移动能耗则可忽略不计。因此在本文的建模中，只考虑任务能耗。

我们定义无人机 $j \in \mathcal{X}_{V}$ 在处理其分配到的任务时所消耗的计算能耗为 $E_j^{comp}$。这部分能耗与其 CPU 的运算功率和任务执行时间成正比。

无人机 $j$ 的计算功率与其CPU频率的三次方成正比：
\begin{equation}
	P_j^{comp} = \kappa_j (F_j)^3
\end{equation}

其中 $\kappa_j$ 是无人机 $j$ 的一个有效电容系数。

无人机 $j$ 在执行任务 $k$ 期间所消耗的计算能耗为：
\begin{equation}
	E_{jk}^{comp} = P_j^{comp} \times T_{jk}^{comp}
\end{equation}

系统在某个时隙内的总计算能耗是所有无人机执行其分配任务所消耗能耗的总和。
对于单个无人机$j$，其总计算能耗是其所分配的所有任务的能耗总和：
\begin{equation}
	E_j^{comp} = \sum_{k \in \mathcal{X}_D} E_{jk}^{comp} 
\end{equation}

其中，$\mathcal{X}_{D}$ 是任务集合。

因此，整个无人机集群的总计算能耗为所有无人机能耗的总和：
\begin{equation}
	E^{sys} = \sum_{j \in \mathcal{X}_{V}} E_j^{comp} = \sum_{j \in \mathcal{X}_{V}} \sum_{k \in \mathcal{X}_D} P_j^{comp} \times T_{jk}^{comp}
\end{equation}

将拓扑优化问题建模为一个基于李雅普诺夫优化的框架。此方法通过引入虚拟队列，将复杂的长期优化问题转化为每个时隙的瞬时优化问题，从而简化求解。

对于随机的任务需求，我们定义虚拟队列$Q_t(t)$用于衡量系统积压的工作量，$Q_e(t)$用于衡量系统的总能耗代价。
\begin{equation}
	Q_t(t)=\max\{Q_t(t-1)+\Psi^{gene}-\Psi^{\text{sys}},0\}
\end{equation}

其中$Q_t(t-1)$表示上一时刻任务队列状态，$\Psi^{gene}$表示产生的任务量，$\Psi_{\text{sys}}$表示解决的任务量。
\begin{equation}
	Q_e(t) = \max\{Q_e(t-1)+E^{\text{sys}}-E^u,0\}
\end{equation}

其中$Q_e(t-1)$表示上一时刻能量队列状态，$E^u$表示无人机平均每时隙消耗能量上界。
定义了李雅普诺夫函数：
\begin{equation}
	L(\Theta(t))=\frac{1}{2}(Q_t^2[t]+Q_e^2[t])
\end{equation}

当$Q_t(t)$任务队列积累较多时，倾向于将无人机部署到用户密度高、需求大的区域，以期减少任务压力、缓解集群压力。

当$Q_e(t)$能量队列积累较多时倾向于减少无人机移动，或将其部署到能耗较低的位置，降低能源消耗、提高续航能力。

针对簇内拓扑位置优化问题，我们将利用经典的李亚普诺夫方法，通过定义李雅普诺夫函数、引入惩罚项，逐步解决问题。将优化目标转化为最小化每个时隙的“漂移-plus-惩罚”项，从而指导无人机的移动决策。我们以无人机位置坐标为决策变量，求解出在下一时刻的优化位置，最终使整个簇达到最优或次优的拓扑状态。

在解决理论集群裂变与汇合问题时，由于这个问题是一个复杂的组合优化问题，难以找到全局最优解，因此启发式算法是一个高效且实用的选择。基于以上的虚拟队列衡量任务进度和无人机能耗，预设动态阈值，通过规则快速推理集群状态。同时考虑用户分布和无人机分布量化集群形态，指导集群进行裂变与汇合决策。

李雅普诺夫优化输出的目标位置即为无人机飞行的目标点，无人机的位置调整则依赖一个安全可靠的轨迹规划算法。。

冲突搜索算法（Conflict-Based Search, CBS）是一种高效的多智能体路径规划算法，它通过迭代地解决冲突来找到最优的路径，为整个集群生成一套全局无碰撞的运动路径。

在CBS的语境中，“障碍”有两种。一种是静态障碍物，如地图上的墙壁、建筑物等。一种是动态障碍物，即其他智能体，智能体彼此之间就是移动的障碍物。CBS通过检测冲突、施加约束、重新规划的迭代循环来解决这些问题。一旦发现两个智能体的路径存在上述冲突，它就会施加约束。

标准CBS不考虑智能体之间的通信连接性。规划出的路径可能使无人机之间距离过远，导致通信链路中断，集群解体。为解决无人机集群特有的通信连接保持问题，我们引入通信感知的冲突约束。

无人机间的传输为空对空传输，通常视为视距传输，适用于自由空间模型。对于任意两个无人机 $u_i$ 和 $u_j$，有信道增益：
\begin{equation}
	h_{i,j}=\frac{1}{\beta \delta^2(v_i,v_j)}\mu_{LoS}|h^{Rice}_{i,j}|^2
\end{equation}

其中 $h^{\text{Rice}}_{i,j} \sim \text{Rice}(\upsilon,\delta)$，是莱斯衰落项。$\beta_0=(\frac{4\pi f}{c})^2$，f是携带频率，c是光速，$\delta(v_i,v_j)$是无人机之间的距离，$\mu_{LoS}$是LoS的额外路径损失参数。

那么其信噪比为：
\begin{equation}
	 SINR^{A2A} = \frac{P^t }{N_0}h_{i,j}
\end{equation}

其中$P^t$ 是传输功率，$N_0$是噪音功率。

数据可达率为
\begin{equation}
	r_{i,j}^{A2A} = B\log(1+SINR_{i,j}^{A2A})
\end{equation}

其中B为链路带宽。
为确保控制信令与任务数据的可靠传输，必须要求该速率不低于维持网络连接所需的最低阈值 $r_{\text{th}}$：
\begin{equation}
	\label{eq:rate_constraint}
	r_{i,j}^{A2A} \geq r_{\text{th}}
\end{equation}

当检测到此类通信冲突时，它将施加相应的通信约束。该约束不再仅仅是禁止某个位置，而是要求智能体 $ξ_i$ 在时间 $t$ 必须位于一个能与其邻居集合 $\Xi$ 保持有效连接的区域内。

基于上述度量，我们形式化地定义一种新的冲突类型——通信冲突。该冲突定义为四元组，$\mathcal{C}_{\text{com}} = \langle \xi_i, \Xi, t, loc \rangle$：

其中$\xi_i$是发生通信中断的智能体；$\Xi$是$\xi_i$ 需要维持通信的邻居智能体集合；$t$是通信中断发生的时刻；$loc$是 $\xi_i$ 在时刻 $t$ 所处的位置。

该冲突表示在当前路径规划下，智能体 $\xi_i$ 在时刻 $t$ 位于位置 $loc$ 时，其与邻居集合 $\Xi$ 中所有智能体的通信链路质量均不满足约束。这意味着当前规划路径可能导致 $\xi_i$ 在时间 $t$ 脱离网络拓扑，需要我们进一步调整。

在集群划分与分配映射子问题中，我们将采用模糊C均值聚类（Fuzzy C-Means Clustering, FCM）算法。FCM是一种经典的聚类算法，它允许一个数据点属于多个簇，具有“模糊”的属性。对于簇网络边缘的无人机，允许其属于多个簇，避免了拓扑网络的剧烈变动，提供了一个过渡的空间。

我们将无人机视为作为一个个数据点，其资源表示为特征向量，包括位置、电量、任务能力等。通过FCM算法，可以实现更加灵活和合理的集群划分，使得每个无人机能够以一定的概率属于不同的子集群，提供更多的可能性。

目标函数为最小化所有数据点到其聚类中心的加权欧氏距离之和。数学表达式如下：
\begin{equation}
	J_m = \sum_{j=1}^{n}\sum_{c=1}^{C} \mu_{jc}^{m} \|x_j - c_c\|^2
\end{equation}

其中$J_m$ 是目标函数，表示所有节点到其聚类中心的加权欧氏距离平方和；$n$ 是服务节点的总数；$C$ 是子集群的数量； $\mu_{jc}$ 是服务节点 $j$ 到子集群中心 $c$ 的隶属度，其值在 $[0, 1]$ 之间； $m$ 是模糊化因子，通常取 $m > 1$； $\|x_j - c_c\|^2$ 是节点 $j$ 的特征向量 $x_j$ 到聚类中心 $c_c$ 的欧氏距离。

为了使 FCM 算法能够根据任务需求和资源状况进行智能划分，我们将每个服务节点的特征向量 $x_j$ 定义为一个多维向量，包含：
$$
x_j = (loc_j, E_j, cap_j)
$$

其中$loc_j = (x_j, y_j, z_j)$，表示节点的三维位置坐标；$E_j$是节点的当前剩余电量；$cap_j = (F_j, M_j)$表示节点的计算能力,如CPU及内存。

通过迭代优化，FCM 算法将计算出每个服务节点对所有潜在子集群的隶属度，从而实现动态的、平滑的集群划分。这种“模糊”划分方式可以有效避免拓扑的剧烈变化，为集群的平稳裂变或汇合提供理论基础。

\subsubsection{高动态服务系统能量均衡路由问题建模}
本研究对网络系统、通信、能量及数据队列的全面建模，
并在此基础上构建了以网络长期稳定运行为目标的优化问题。
各参数定义如表\ref{tab:variable}所示，在这个问题中我们将高动态边缘移动设备特化为无人机进行表述。
\begin{table}[h]
\centering
\caption{使用的参数符号表}
\begin{tabular}{|l|l|}
\hline
\textbf{符号} & \textbf{描述} \\
\hline
$N$ & 集群中UAV的总数量 \\
\hline
$D$ & 数据中心的地理位置 \\
\hline
$q_i(t)$ & UAV$i$在时间$t$的地理位置和速度信息 \\
\hline
$d_{ij}(t)$ & 节点$i$和$j$在时间$t$的欧几里得距离 \\
\hline
$P_{ij}^{trans}(t)$ & UAV$i$在时间$t$对UAV$j$的传输功耗 \\
\hline
$P_{max}^{trans}$ & UAV的最大传输功耗 \\
\hline
$r_{ij}(t)$ & UAV$i$在时间$t$对UAV$j$的数据传输速率 \\
\hline
$x_{i,j}(t)$ & UAV$i$在时间$t$是否对UAV$j$转发数据 \\
\hline
$E_{i}(t)$ & UAV$i$在时间$t$剩余能量 \\
\hline
$E^{max}$ & UAV通信能量预算 \\
\hline
$Q_{i}(t)$ & UAV$i$在时间$t$数据队列堆积 \\
\hline
$Q^{max}$ & UAV最大数据队列堆积 \\
\hline
$T$ & 一次任务执行的总时间 \\
\hline
$F_{i}(t)$ & UAV$i$在时间$t$的虚拟能量队列堆积 \\
\hline
\end{tabular}
\label{tab:variable}
\end{table}

假设有$\mathcal{N}$个无人机构成的无人机群,描述为$U=\{1,2,...,\mathcal{N}\}$，
每架无人机的坐标信息为 $q_{i}=(x_{i}^{uav},y_{i}^{uav})$，
速度信息为 $v_{i}=\{v_{xi}^{uav},v_{yi}^{uav}\}$，这些信息均是关于时间的函数 $q(t)$。

假设无人机间没有障碍物，通信以视距链路为主，信道增益 $h_{ij}$ (相当于接收功率的衰减比例)计算为
\begin{equation}
    h_{ij}=\frac{\rho_{0}}{d_{ij}^{2}}
\end{equation}
其中$\rho_{0}$为参考距离下的信道功率，$d_{ij}$为节点间距离$\rho_{0}$表示为$\rho_{0}=\frac{G_{i}G_{j}\lambda^{2}}{(4\pi)^{2}}$，
$G$ 为两节点的天线增益(固定硬件参数)，$\lambda$为信号波长，由载波频率决定，通常固定，因此可视为常数。
信噪比（SINR）用于量化有用信号强度相对于干扰信号和背景噪声的综合影响，则节点 $i,j$ 间的信噪比定义为
\begin{equation}
    SINR_{ij}=\frac{h_{ij}P_{ij}^{trans}(t)}{I_{ij}(t)+\sigma^{2}}
\end{equation}
其中 $P_{ij}^{trans}(t)$ 为发送功率，$I_{ij}(t)$ 为干扰功率，$\sigma^{2}$ 为噪声功率,噪声功率视为固定值。
干扰功率 $I_{ij}(t)$ 表示节点 $j$ 在接收来自节点 $i$ 的信号时，网络中其他所有节点对其产生的干扰功率总和，建模为
\begin{equation}
    \bar I_{ij}=\sum_{k \neq i} \sum_{l \neq j} {h_{kj}P_{kl}^{trans}\cdot x_{kl}(t)}
\end{equation}
数据在无人机之间的传输速率 $r_{ij}(t)$ 可由香农容量模型进行计算,则可表示为
\begin{equation}
    r_{ij}(t)=B~\log_{2}(1+SINR_{ij}(t))
\label{eq:rij}
\end{equation}
其中B为带宽,表述为信号可用的最高频率和最低频率之差，可以为固定值。
无人机在时刻 $t$ 的总功耗 $P_{i}(t)$ 仅由传输功耗构成，每个无人机的电量变化可以表示为:
\begin{equation}
    E_{i}(t+1)=E_{i}(t)-P_{i}(t)\cdot\Delta t
\end{equation}
约束 $E^{min}\le E_{i}(t)\le E^{max}$

每个无人机的队列长度变化表示为:
\begin{equation}
    Q_{i}(t+1)=\max\{Q_{i}(t)-\mu_{i}(t)\cdot\Delta t,0\}+\nu_{i}(t)\cdot\Delta t
\end{equation}
其中 $\Delta t$ 是决策时间间隔，$\mu_{i}(t)=\sum_{j}x_{ij}(t)r_{ij}(t)$是总发出数据速率，
$\nu_{i}(t)=\lambda_{i}(t)+\sum_{k}x_{ki}(t)r_{ki}(t)$总到达数据速率，约束 $Q_{i}(t)\le Q^{max}$。

因此优化目标可以定义为：
\begin{align}
\text{Maximize} & \quad \frac{1}{T}\sum_{t=0}^{T} \sum_{i=1}^{N} \mu_i(t) \\
\text{s.t.} \quad 
&0 \le P_{ij}^{\text{trans}}(t) \le P_{\text{max}}^{\text{trans}}, \quad \forall i,j,t \label{C1}\\
& \sum_{j \in \mathcal{N} \cup \{D\}, j \neq i} x_{ij}(t) \le 1, \quad \forall i,t \label{C2}\\
& 0 \le Q_i(t) \le Q^{\text{max}}, \quad \forall i,t \label{C3}
\end{align}

其中约束\ref{C1}确保对于所有时隙，传输功率都在限制范围内，每个无人机最多转发到一个邻居（约束\ref{C2}），数据队列不会溢出（约束\ref{C3}）。

为了处理长期能量预算约束，引入虚拟队列 $F_{i}(t)$ 来追踪累积的超额能量消耗：
\begin{equation}
    F_{i}(t+1)=\max\{F_{i}(t)+P_{i}(t)\cdot\Delta t-P_{avg}\cdot\Delta t,0\}
\end{equation}
其中 $P_{i}(t)$ 是无人机$i$在时间$t$的总功率消耗，$P_{avg}$ 是允许的平均功率预算，可根据任务持续时间和电池容量设定为$P_{avg}=\frac{E^{max}}{T}$。


定义李雅普诺夫函数,结合队列稳定性和能量平衡
\begin{equation}
    L_{i}(t)=\frac{1}{2}[Q_{i}(t)]^{2}+\frac{\beta}{2}[\Phi_{i}(t)]^{2}+\frac{\gamma}{2}[F_{i}(t)]^{2}
\end{equation}
其中$\Phi_{i}(t)=E^{max}-E_i(t)$表示当前能量损失。一步条件李雅普诺夫漂移可表示为$\Delta L_{i}(t) = \mathbb{E}[L(t+1)-L(t)|Q(t),\Phi(t),F(t)]$
测量系统不稳定性的预期变化。漂移加惩罚框架用于通过寻求最小化漂移的上界减去表示所需总网络吞吐量的$V$加权惩罚项来联合优化稳定性和网络吞吐量。
这导致漂移加惩罚表达式的以下上界：
\begin{equation}
\begin{split}
\Delta L_{i}(t) - V\mathbb{E}[\mu_i(t)|Q(t)] \le C &+ \mathbb{E}[Q_{i}(t)\Delta t(\nu_{i}(t)-\mu_{i}(t))|Q(t)] \\
&+ \mathbb{E}[(\beta\Phi_{i}(t)+\gamma F_{i}(t))\Delta tP_{i}(t)|\Phi(t),F(t)] \\
&- \mathbb{E}[\gamma F_{i}(t)\Delta tP_{avg}|F(t)] \\
&- V\mathbb{E}[\mu_i(t)|Q(t)]
\end{split}
\end{equation}
其中$C$是一个常数，每个UAV$i$的每个时隙目标是最小化这个界限，这是通过只关注期望范围内的条件最小化来实现的：
\begin{equation}
\min_{\mathbf{x}_i(t),\mathbf{P}_i^{trans}(t)} \bigg\{ Q_i(t)\Delta t\nu_i(t) - D_i(t)\mu_i(t) + C_i(t)\Delta tP_i(t) \bigg\}
\end{equation}
其中$\mathbf{x}_i(t)=\{x_{ij}(t)\}_{\forall j}$和$\mathbf{P}_i^{trans}(t)=\{P_{ij}^{trans}(t)\}_{\forall j}$是UAV$i$在时间$t$的路由决策和功率控制，$C_{i}(t)=\beta\Phi_{i}(t)+\gamma F_{i}(t)$, $D_{i}(t)=Q_{i}(t)\Delta t+V$。

\subsubsection{基于启发式的分布式高动态边缘服务设备路由决策和转发功率选择}
\begin{algorithm}[!ht]
    \caption{LEADR分布式路由和功率控制}
    \label{alg:LEADR}
    \renewcommand{\algorithmicrequire}{\textbf{Input:}}
    \renewcommand{\algorithmicensure}{\textbf{Output:}}
    \begin{algorithmic}[1]
        \REQUIRE 本地状态$Q_i(t)$，$E_i(t)$，$F_i(t)$，邻居状态来自接收到的HELLO消息。
        \ENSURE 控制动作$x_{ij}(t)$和$P_{ij}^{trans}(t)$
        \FOR{邻居$j \in N_i(t) \cup D$}
            \STATE 计算最佳传输功率$P_{ij}^{\text{trans,opt}}(t)$（公式\ref{eq:pij}）。
            \STATE 计算相应的数据速率$r_{ij}^{\text{opt}}(t)$（公式\ref{eq:rij}）。
            \STATE 计算复合启发式权重$W_{ij}^{\text{E,G}}(t)$（公式\ref{eq:wij}）。
        \ENDFOR
        \STATE 选择最佳下一跳邻居：$j^{*}(t)=\arg\max_{j\in N_{i}(t)\cup\{D\}}W_{ij}^{E,G}(t)$。
        \STATE 设置路由决策 $x_{ij}^{*}(t)$。
        \STATE 将数据传输给邻居 $j^{*}(t)$ 以功率 $P_{ij}^{trans,opt}(t)$。
        \STATE 更新状态 $Q_i(t+1)$,$E_i(t+1)$,$F_i(t+1)$ 用于下一个时隙。
        \RETURN 转发决策 $x_{ij}(t)$ 和转发功率 $P_{ij}^{trans,opt}(t)$。
    \end{algorithmic}
\end{algorithm}
算法\ref{alg:LEADR}显示了我们提出的算法的过程。
对于每个潜在的下一跳邻居$j$，UAV$i$首先计算最佳发射功率$P_{ij}^{trans，opt}(t)$。
这是通过隔离目标函数中与单个链路$(i，j)$相关的项并解决以下子问题来实现的：
\begin{equation}
    \min_{P_{ij}^{trans}(t)}f(P_{ij}^{trans}(t))=-D_{i}(t)\cdot B~\log_{2}\left(1+\frac{h_{ij}P_{ij}^{trans}(t)}{N_{ij}}\right)+C_{i}(t)\cdot\Delta t\cdot P_{ij}^{trans}(t)
\end{equation}
其中$N_{ij}=I_{ij}(t)+\sigma^{2}$。此函数是$P_{ij}^{trans}(t)$的凸函数。
通过将其一阶导数设置为零，产生最佳幂$P_{ij}^{trans*}(t)$。考虑到物理约束，最终传输功率定义为：
\begin{equation}
    P_{ij}^{trans,opt}(t)=\min\{\max\{0,P_{ij}^{trans*}(t)\},P_{max}^{trans}\}
\label{eq:pij}
\end{equation}
在计算出每个潜在邻居的最佳功率后，UAV$i$必须选择单个最佳下一跳。
该决策由综合链路权重指导，该权重结合了实用的路由启发式方法。
首先，通过将最佳功率和结果速率插入目标函数来计算基本权重$W_ij(t)$。该权重代表选择链路$(i,j)$的收益：
\begin{equation}
    W_{ij}(t)=D_{i}(t)\cdot r_{ij}^{opt}(t)-C_{i}(t)\cdot\Delta t\cdot P_{ij}^{trans,opt}(t)
\end{equation}
为了防止路由循环和转发到能量不稳定的节点，然后通过结合两个关键启发式来修改这个基本权重。
第一个启发式通过其归一化能量$\tilde{E}_{j}(t)=\frac{E_{j}(t)}{E^{max}}$来评估潜在下一跳$j$的能量健康状况，这不鼓励转发到电池电量低的无人机。
第二个启发式通过计算相对距离参数$D_{ij}(t)=\frac{d_{iD}(t)-d_{jD}(t)}{d_{iD}(t)}$，并且只允许转发到$D_{ij}(t)>0$的节点：
\begin{equation}
    W_{ij}^{E,G}(t)=\begin{cases}W_{ij}(t)\cdot\tilde{E}_{j}(t)& \text{if } D_{ij}(t)>0 \text{ and } j\ne D \\ W_{ij}(t)& \text{if } j=D \text{ and } D_{ij}(t)>0  \\ -\infty& \text{otherwise}\end{cases}
\label{eq:wij}
\end{equation}
然后，UAV选择最大化此复合权重的邻居$j^{*}(t)$：
\begin{equation}
    j^{*}(t)=\arg\max_{j\in N_{i}(t)\cup\{D\}}W_{ij}^{E,G}(t)
\end{equation}
\begin{equation}
    x_{ij}(t)=\begin{cases}1& \text{if } j=j^{*}(t) \text{ and } W_{ij^{*}}^{E,G}(t)>-\infty \\ 0& \text{otherwise}\end{cases}
\end{equation}
\subsection{取得的成果}
目前在整个课题的相关研究中取得的成果：
\begin{enumerate}
  \item \textbf{论文在投：“LEADR：一种基于Lyapunov的连续无人机通信服务能量感知去中心化路由策略”}。本研究通过将任务长时间的最佳化问题转化为每架无人机的一系列低复杂性、每时隙决策来解决无人机路由能耗和吞吐量不平衡的挑战。通过将李雅普诺夫漂移加惩罚框架与基于地理和能量的启发式方法相结合，LEADR能够实现主动能量平衡管理，同时保持高网络吞吐量。
  \item \textbf{专利受理：“面向微服务架构无人机网络的联合路由与功率控制方法”}。该发明包含以下步骤，步骤1：无人机系统状态建模；步骤2：基于李雅普诺夫优化的联合控制问题构建函数；步骤3：分布式路由与功率联合决策；步骤4：一种基于链路寿命预测的自适应信息交互机制；步骤5：无人机网络性能仿真与综合评估。通过比较不同场景下网络总吞吐量、网络中所有无人机剩余电量的均衡程度、以及网络能够持续工作的时间等关键性能指标，验证了本发明方法的有效性和优越性。
\end{enumerate}
\subsection{进度安排}
本项目共分为7部分，详细时间安排见表\ref{tab:Schedule}.
\begin{table}[ht]
    \vspace{0.5em}\centering\wuhao
    \caption{\label{tab:Schedule} 进度安排}
    \begin{tabular}{>{\centering\arraybackslash}p{0.2\textwidth}>{\centering\arraybackslash}p{0.55\textwidth}>{\centering\arraybackslash}p{0.2\textwidth}}
        \toprule[1.5pt]
        \textbf{时间} & \textbf{进度安排} \\
        \midrule[1pt]
        2025/9-2025/10 & 相关领域文献调研，完善并最终确定三个研究问题的详细技术方案\\
        2025/11-2026/2 & 完成物理拓扑优化方法的算法设计与实现，完成与基线算法的对比实验与性能验证。\\
        2026/3-2026/6 & 完成能量均衡路由策略的算法设计与实现，完成与基线算法的对比实验与性能验证。 \\
        2026/7-2026/10 & 完成基于深度强化学习的联合治理框架设计与实现，进行消融实验，验证联合优化的协同增益 \\
        2026/11-2026/12 & 对全部研究工作进行系统性总结、理论分析与实验结果整理\\
        2026/12 & 根据导师指导意见，修改并完善学位论文，完成论文的审查与提交\\
        2027/1 & 准备并完成硕士学位论文的最终答辩\\
        \bottomrule[1.5pt]
    \end{tabular}
\end{table}
\section{为完成课题已具备和所需的条件和经费}
目前已基本掌握本课题需要使用的具体技术，
对高动态网络中的拓扑控制、移动自组织网络路由协议、
李雅普诺夫优化理论以及深度强化学习（特别是多智能体深度强化学习）
等核心理论与技术方法有了系统性的了解和掌握。
后续实验所需要的计算资源和设备都已具备，满足研究所需条件。
\section{预计研究过程中可能遇到的困难和问题，以及解决措施}
问题1：深度强化学习模型在高动态边缘场景中的收敛性与泛化能力挑战。

解决方案：采用多目标奖励机制，不仅奖励数据成功转发，还引入对能耗均衡度、网络连通性和簇结构稳定性的量化奖励/惩罚项。

问题2：李雅普诺夫优化理论假设与极端网络动态性的适配问题。

解决方案：在算法设计中，除了遵循李雅普诺夫框架推导出的主决策逻辑外，
将引入一套启发式的“安全保护”机制。
例如，当节点的数据队列长度或能量赤字超过预设的紧急阈值时，
将触发一个备用的、更保守的路由策略（如优先保障数据传出的贪心路由），
以牺牲短期最优性为代价来确保系统的稳定性。
\section{主要参考文献}
\bibliographystyle{hithesis}
\bibliography{reference}

% Local Variables:
% TeX-master: "../mainart"
% TeX-engine: xetex
% End: